\documentclass[a4paper,12pt,twoside]{report}
\usepackage{polyglossia}
\setdefaultlanguage{croatian}
\setotherlanguage[variant=ancient]{greek}

%\usepackage{fontspec}
%\usepackage{verse}
\defaultfontfeatures{Ligatures=TeX}

\usepackage{import}
\usepackage[small,sf,bf]{titlesec}
\usepackage{tabto}
\usepackage{ulem}
\usepackage{hyperref}
\usepackage{enumitem}
\usepackage{dtk-logos}

\usepackage{graphicx}
\graphicspath{ {./slike/} }

\usepackage[all]{nowidow}

\usepackage{fancyhdr}
\renewcommand{\chaptermark}[1]{\markboth{#1}{}}
\renewcommand{\sectionmark}[1]{\markright{#1}}
\pagestyle{fancy}
\fancyhf{}
\fancyhead[LE,RO]{\thepage}
\fancyhead[RE]{\itshape\nouppercase{Grčki rečenični period}}
\fancyhead[LO]{\nouppercase{\leftmark}}
\renewcommand{\headrulewidth}{0pt}

\usepackage{titling}
\newcommand{\subtitle}[1]{%
  \posttitle{%
    \par\end{center}
    \begin{center}\large#1\end{center}
    \vskip0.5em}%
}
 
\setmainfont{Old Standard TT}
\setsansfont{Old Standard TT}


\hyphenation{δυσ-σέ-βει-αν βού-λεσ-θαί κα-τη-γο-ρού-σης τα-χέ-ως πε-πλημ-μέ-λη-κε νε-α-νί-ας αὐ-τῇ ἀ-φε-λό-με-νος ἀ-πή-γε-το ἐ-πέ-πληξ-άς ἠ-γό-μην ἤ-νεγ-κεν νο-μί-ζει ἄν-θρω-πον παν-το-δα-πά Αἰ-θί-οψ-ιν ἀν-έρ-χε-ται φρον-τί-δων αὐ-τὸς δι-ῃ-ρη-μέ-νος κλέπ-τον-τας ἑ-κα-τόμ-βῃ μέ-γισ-τον κιν-δύ-νου}

\hyphenation{Ελ-λή-νων ξυν-έ-μει-νεν ἐ-πι-όντ-ων ἀνα-σκευ-α-σά-με-νοι ἀ-πω-σά-με-νοι Λα-κε-δαι-μό-νι-οι Ελ-λη-νες δι-ε-φά-νη}

\hyphenation{εὐ-δο-κι-μή-σας ἐν-ταῦ-θα βα-σι-λεύ-εις ἐ-χρή-σα-το συλ-ληφ-θέν-τος}

\hyphenation{θε-ρά-πον-τες ἡ-γοῦ-μαι}

\hyphenation{ἐπι-φα-νέσ-τε-ρον το-σοῦ-τον ἔ-χον-τας συγ-γιγ-νο-μέ-νους μᾶλ-λον με-γίσ-την}

\hyphenation{ἐπ-έσ-κηπ-τε Θε-μισ-το-κλῆς}

\hyphenation{κοι-νω-νοῦ-σιν δη-λῶ-σαι δι-α-λε-γο-μέ-νους πλη-σι-ά-ζον-τας πα-ρα-λι-πεῖν}

\hyphenation{καρ-πῶν ὀ-νο-μα-ζό-με-νον τηκ-τὰ ὅ-σα}

\hyphenation{με-μά-θη-κας Μά-λισ-τα πο-λυ-μα-θής}

\hyphenation{δι-α-βε-βλη-μέ-νος Α-ρι-στό-βου-λος Σω-κρά-τους Δı-α-τ
ρι-βαί ἐξ-ελ-κύ-σαι Ἐγ-χει-ρί-δι-ον}

\hyphenation{δά-μα-λιν δύ-να-μαι}

\hyphenation{πυ-θο-μέ-νου ἀ-πο-κρί-νε-ται τρα-πέ-ζαις ἀ-πέσ-τει-λαν γρά-ψαν-τος με-τα-τί-θη-σι με-γά-λου Τισ-σα-φέρ-νῃ προσ-ε-πι-μετ-ρῆ-σαι Α-θη-ναί-οις Τισ-σα-φέρ-νην Αλ-κι-βι-ά-δην πα-τρί-δα}

\hyphenation{ἡ-γοῦν-ται ἀ-πο-ρω-τά-των πό-λεις Ras-pra-va Ἀ-πο-λο-γί-α Λα-κε-δαı-μο-νί-ων Κυ-νη-γε-τι-κός}

\hyphenation{πλε-όν-των γε-ωρ-γὸς ἕ-τε-ρον}

\hyphenation{κα-θεῖ-ναι προ-ελ-θοῦ-σαν ἀ-πο-δοῦ-ναι ἀ-πο-φαί-νον-τος ἀ-πεσ-τά-λη ἅ-παν-τας τρί-πο-δα τρί-πο-δος κα-θι-ερώ-θη}

\hyphenation{χρεί-αν παν-τὸς ἀν-επί-σκεπ-τον ἀ-πο-θα-νόν-των πα-ρα-σκευ-ά-ζον-τας φρον-τί-ζον-τας πολ-λοὺς κτω-μέ-νους ἐ-λατ-τοῦσ-θαι ἀ-με-λοῦν-τας ἀ-θε-ρά-πευ-τον κτη-μά-των πει-ρω-μέ-νους ὀ-λι-γω-ροῦν-τας δε-ο-μέ-νων ἐ-ῶν-τας ἔ-μοι-γε}

\hyphenation{αὐ-το-κρά-το-ρας Λά-μα-χον ποι-ή-σαν-τες ἑξή-κον-τα Σε-λι-νουν-τί-ους χρη-μά-των μισ-θόν Νι-κη-ρά-του πε-ρι-γίγ-νη-ται}

\def\secondpage{\clearpage\null\vfill
\pagestyle{empty}
\begin{minipage}[b]{0.9\textwidth}
\begin{center}
\textit{Izdavač}\\
Filozofski fakultet Sveučilišta u Zagrebu\\
FF-Press\\
Ivana Lučića 3, Zagreb

\bigskip

\textit{Za izdavača}\\
prof.\ dr.\ Vesna Vlahović Štetić

\bigskip

\textit{Urednici}\\
Nina Čengić\\
Neven Jovanović

\bigskip

\textit{Recenzenti}\\
NN\\
NN

\bigskip

\textit{Lektor}\\
Nina Čengić

\bigskip

Korekturu i računalni slog obavili autori

\bigskip

%ISBN 978-953-175-817-8

\bigskip

Ovo izdanje elektroničke knjige zaključeno je u rujnu 2019.

%\bigskip

%\includegraphics[scale=0.20]{ffpress-cb.png}


\end{center}
\footnotesize\raggedright
\setlength{\parskip}{0.5\baselineskip}
\noindent Ovo djelo je ustupljeno pod Creative Commons licencom Imenovanje 3.0 nelokalizirana licenca. Da biste vidjeli primjerak te licence, posjetite \url{http://creativecommons.org/licenses/by/3.0/} ili pošaljite pismo na Creative Commons, PO Box 1866, Mountain View, CA 94042, SAD. 
\end{minipage}
\vspace*{3\baselineskip}
%\cleardoublepage
\rfoot{\thepage}}

\makeatletter
\g@addto@macro{\secondpage}{\maketitle}
\makeatother

\begin{document}

\title{Čitanka za kolegij Povijest grčkoga jezika: Grčki rečenični period klasičnoga i poklasičnoga razdoblja}
%\subtitle{Vježbe 41–60}
\author{Nina Čengić, Neven Jovanović}
\date{}
\secondpage


\cleardoublepage
\setcounter{page}{2}
\thispagestyle{empty}


%\frontmatter


\chapter*{Predgovor}
\label{chap:predgovor}
\addcontentsline{toc}{chapter}{\nameref{chap:predgovor}}

\section*{O ovoj čitanci}

Ovaj izbor komentiranih tekstova namijenjen je upoznavanju strukture grčke rečenice, i dio je literature obaveznog kolegija \textit{Povijest grčkoga jezika: Grčki rečenični period klasičnoga i poklasičnoga razdoblja}. Izbor uključuje autore poput Herodota i Ezopa, Platona, Ksenofonta i Tukidida, Aristotela, Lizije, Izokrata i Demostena, pa sve do Polibija, Marka Aurelija, Filostrata, Ahileja Tacija i Plotina; zastupljeni su historiografija i filozofska proza, govorništvo, pripovijesti, basne.

Za uspješniji samostalni rad kratkim je odlomcima dodan uvod, u kojem su predstavljeni autor (kada se susreće prvi put) i tekst iz kojeg odlomak potječe.

Čitanka je priređena računalnim programima za slaganje teksta \LaTeX\ i \XeLaTeX. Izvorni kod dostupan je u repozitoriju Bitbucket, na URL adresi \href{https://bitbucket.org/nevenjovanovic/hellenismos-g}{bitbucket.org/nevenjovanovic/hellenismos-g}.

\medskip

U Zagrebu, rujan 2019.

\newpage

\tableofcontents

\newpage

%\mainmatter
\chapter[Αἰσώπειος μῦθος 216]{\textgreek[variant=ancient]{Αἰσώπειος μῦθος} 216}

\import{poglavlja/}{pogl01-ezop-216.tex}

\chapter[Πλάτωνος Μενέξενος]{\textgreek[variant=ancient]{Πλάτωνος Μενέξενος} 237c}

\import{poglavlja/}{pogl02-platon-meneksen237c.tex}

\chapter[Λουκιανοῦ Δὶς κατηγορούμενος ]{\textgreek[variant=ancient]{Λουκιανοῦ Δὶς κατηγορούμενος } 2, 9}

\import{poglavlja/}{pogl03-lukijan-dvaput-2-9.tex}


\chapter[Θουκυδίδου Ἱστορίαι Α]{\textgreek[variant=ancient]{Θουκυδίδου Ἱστορίαι Α} 18, 1}

\import{poglavlja/}{pogl04-tukidid-1-18-1.tex}


\chapter[Αἰσώπειος μῦθος 83]{\textgreek[variant=ancient]{Αἰσώπειος μῦθος} 83}

\import{poglavlja/}{pogl05-ezop-83.tex}

\chapter[Πλάτωνος Φαίδων]{\textgreek[variant=ancient]{Πλάτωνος Φαίδων} 85a}

\import{poglavlja/}{pogl06-platon-fedon85a.tex}

\chapter[Ἰσοκράτους Βούσιρις]{\textgreek[variant=ancient]{Ἰσοκράτους Βούσιρις} 28, 1}

\import{poglavlja/}{pogl07-izokrat-buzirid28.tex}

\chapter[Θουκυδίδου Ἱστορίαι Ϛ]{\textgreek[variant=ancient]{Θουκυδίδου Ἱστορίαι Ϛ} 30, 1}

\import{poglavlja/}{pogl08-tukidid-6-30-1.tex}

\chapter[Λυσίου Κατὰ Ἀγοράτου]{\textgreek[variant=ancient]{Λυσίου Κατὰ Ἀγοράτου} 39}

\import{poglavlja/}{pogl09-lizija-agorat-39-1.tex}

\chapter[Ἰσοκράτους Εὐαγόρας]{\textgreek[variant=ancient]{Ἰσοκράτους Εὐαγόρας} 9}

\import{poglavlja/}{pogl10-izokrat-euagora-9.tex}

\chapter[Πλάτωνος Κριτίας]{\textgreek[variant=ancient]{Πλάτωνος Κριτίας} 114e}

\import{poglavlja/}{pogl11-platon-kritija-114e.tex}

\chapter[Ἰσοκράτους Πρὸς Δημόνικον]{\textgreek[variant=ancient]{Ἰσοκράτους Πρὸς Δημόνικον} 16}

\import{poglavlja/}{pogl12-izokrat-demoniku-16.tex}

\chapter[Ἀρριανοῦ Ἀναβάσεως Ἀλεξάνδρου Β]{\textgreek[variant=ancient]{Ἀρριανοῦ Ἀναβάσεως Ἀλεξάνδρου Β,} 3, 6}

\import{poglavlja/}{pogl13-arijan-anabaza-2-3-6.tex}

\chapter[Λουκιανοῦ Θεῶν διάλογοι]{\textgreek[variant=ancient]{Λουκιανοῦ Θεῶν διάλογοι Κ} 7}

\import{poglavlja/}{pogl14-lukijan-razgovoribogova-20-7.tex}

\chapter[Πλουτάρχου Περὶ ἀδολεσχίας]{\textgreek[variant=ancient]{Πλουτάρχου Περὶ ἀδολεσχίας} 513a}

\import{poglavlja/}{pogl15-plutarh-obrbljavosti-513a.tex}

\chapter[Ξενοφῶντος Συμπόσιον Δ]{\textgreek[variant=ancient]{Ξενοφῶντος Συμπόσιον Δ} 34, 4}

\import{poglavlja/}{pogl16-ksenofont-simpozij-4-34-4.tex}

\chapter[Λουκιανοῦ Ἰκαρομένιππος]{\textgreek[variant=ancient]{Λουκιανοῦ Ἰκαρομένιππος \\ἢ Ὑπερνέφελος} 25, 3}

\import{poglavlja/}{pogl17-lukijan-ikaromenip-25-3.tex}

\chapter[Πλουτάρχου Σόλων]{\textgreek[variant=ancient]{Πλουτάρχου Βίοι παράλληλοι, Σόλων} 4, 3}

\import{poglavlja/}{pogl18-plutarh-solon-4-3.tex}

\chapter[Ξενοφῶντος Ἀπομνημονεύματα Β]{\textgreek[variant=ancient]{Ξενοφῶντος Ἀπομνημονεύματα Β} 4, 1}

\import{poglavlja/}{pogl19-ksenofont-uspomene-2-4.tex}



\chapter[Θουκυδίδου Ἱστορίαι Ϛ]{\textgreek[variant=ancient]{Θουκυδίδου Ἱστορίαι Ϛ} 8, 1}

\import{poglavlja/}{pogl20-tukidid-6-8-1.tex}

% 20

\chapter[Ἰσοκράτους Ἑλένης ἐγκώμιον]{\textgreek[variant=ancient]{Ἰσοκράτους Ἑλένης ἐγκώμιον} 35}

\import{poglavlja/}{pogl01-izokrat-helena86.tex}

\chapter[Λυσίου Ὑπὲρ τοῦ ἀδυνάτου]{\textgreek[variant=ancient]{Λυσίου Ὑπὲρ τοῦ ἀδυνάτου} 6, 2}

\import{poglavlja/}{pogl02-lizija-invalid-6-2.tex}

\chapter[Δημοσθένους Ὀλυνθιακός Α]{\textgreek[variant=ancient]{Δημοσθένους Ὀλυνθιακός Α,} 15}

\import{poglavlja/}{pogl03-demosten-olint1-15.tex}

\chapter[Ψευδο-Ξενοφῶντος Ἀθηναίων πολιτεία Β]{\textgreek[variant=ancient]{Ψευδο-Ξενοφῶντος Ἀθηναίων πολιτεία \\Β} 14, 12}

\import{poglavlja/}{pogl04-ps-ksen-2-14-2.tex}

\chapter[Ἀριστοτέλους Ῥητορική]{\textgreek[variant=ancient]{Ἀριστοτέλους Ῥητορική} 1393b}

\import{poglavlja/}{pogl05-aristotel-retorika1393b.tex}

\chapter[Αἰσώπειος μῦθος 1]{\textgreek[variant=ancient]{Αἰσώπειος μῦθος} 1}

\import{poglavlja/}{pogl06-ezop-basne1.tex}

\chapter[Ἀριστοτέλους Ῥητορική]{\textgreek[variant=ancient]{Ἀριστοτέλους Ῥητορική} 1358a}

\import{poglavlja/}{pogl07-aristotel-retorika1358a.tex}

\chapter[Γοργίου Ἑλένης ἐγκώμιον]{\textgreek[variant=ancient]{Γοργίου Ἑλένης ἐγκώμιον} 11, 13}

\import{poglavlja/}{pogl08-gorgija-pohvala-11.tex}

\chapter[Πολυβίου Ἱστορίων Γ]{\textgreek[variant=ancient]{Πολυβίου Ἱστορίων Γ,} 111, 3}

\import{poglavlja/}{pogl09-polibije-povijest-3.tex}

\chapter[Πλουτάρχου Θεμιστοκλῆς]{\textgreek[variant=ancient]{Πλουτάρχου Βίοι παράλληλοι, \\Θεμιστοκλῆς} 2}

\import{poglavlja/}{pogl10-plutarh-temistoklo-2.tex}

\chapter[Λουκιανοῦ Κυνικός]{\textgreek[variant=ancient]{Λουκιανοῦ Κυνικός} 18}

\import{poglavlja/}{pogl11-lukijan-cinik-18.tex}

\chapter[Πλουτάρχου Δημοσθένης]{\textgreek[variant=ancient]{Πλουτάρχου Βίοι παράλληλοι, \\Δημοσθένης} 11}

\import{poglavlja/}{pogl12-plutarh-demosten-11.tex}

\chapter[Ξενοφῶντος Ἀπομνημονεύματα Γ]{\textgreek[variant=ancient]{Ξενοφῶντος Ἀπομνημονεύματα Γ} 12, 4}

\import{poglavlja/}{pogl13-ksenofont-uspomene-3-12.tex}

\chapter[Ἰσοκράτους Ἑλένης ἐγκώμιον]{\textgreek[variant=ancient]{Ἰσοκράτους Ἑλένης ἐγκώμιον} 18}

\import{poglavlja/}{pogl14-izokrat-pohvalahelene-18.tex}

\chapter[Δημοσθένους Ὀλυνθιακός Β]{\textgreek[variant=ancient]{Δημοσθένους Ὀλυνθιακός Β,} 18}

\import{poglavlja/}{pogl15-demosten-olintski2-18.tex}

\chapter[Ἀριστοτέλους Ῥητορική]{\textgreek[variant=ancient]{Ἀριστοτέλους Ῥητορική} 1394a}

\import{poglavlja/}{pogl16-aristotel-retorika-1394a.tex}

\chapter[Αἰσώπειος μῦθος 35]{\textgreek[variant=ancient]{Αἰσώπειος μῦθος} 35}

\import{poglavlja/}{pogl17-ezop-35.tex}

\chapter[Γοργίου Ἑλένης ἐγκώμιον 11, 41]{\textgreek[variant=ancient]{Γοργίου Ἑλένης ἐγκώμιον} 11, 41}

\import{poglavlja/}{pogl18-gorgija-pohvalahelene-11-41.tex}

\chapter[Πλουτάρχου Ἀγεσίλαος]{\textgreek[variant=ancient]{Πλουτάρχου Βίοι παράλληλοι, \\Ἀγεσίλαος} 6}

\import{poglavlja/}{pogl19-plutarh-agesilaj-6.tex}

\chapter[Πλουτάρχου Περὶ παίδων ἀγωγῆς]{\textgreek[variant=ancient]{Πλουτάρχου Περὶ παίδων ἀγωγῆς} 9}

\import{poglavlja/}{pogl20-plutarh-liberiseducandis-2c.tex}

%40
\chapter[Διοδώρου Βιβλιοθήκης ἱστορικής Α]{\textgreek[variant=ancient]{Διοδώρου Σικελιώτου \\Βιβλιοθήκης ἱστορικής Α,} 8, 5}

\import{poglavlja/}{pogl01-diodor-knjiznica.tex}

\chapter[Ἀριστοτέλους Ῥητορική]{\textgreek[variant=ancient]{Ἀριστοτέλους Ῥητορική} 1393a}

\import{poglavlja/}{pogl02-aristotel-retorika-1393a.tex}

\chapter[Θουκυδίδου Ἱστορίαι Α 2]{\textgreek[variant=ancient]{Θουκυδίδου Ἱστορίαι Α} 2}

\import{poglavlja/}{pogl03-tukidid-1-2.tex}

\chapter[Ἡροδότου Ἱστορίαι Α 141]{\textgreek[variant=ancient]{Ἡροδότου Ἱστορίαι Α} 141}

\import{poglavlja/}{pogl04-herodot-1-41.tex}

\chapter[Αἰσώπειος μῦθος 9]{\textgreek[variant=ancient]{Αἰσώπειος μῦθος} 9}

\import{poglavlja/}{pogl05-ezop-9.tex}

\chapter[Ἡλιοδώρου Αἰϑιοπικὰ Α]{\textgreek[variant=ancient]{Ἡλιοδώρου Αἰϑιοπικὰ Α} 2}

\import{poglavlja/}{pogl06-heliodor-aethiop-1.tex}

\chapter[Θουκυδίδου Ἱστορίαι Β]{\textgreek[variant=ancient]{Θουκυδίδου Ἱστορίαι Β} 65}

\import{poglavlja/}{pogl07-tukidid-2-65.tex}

\chapter[Ἰσοκράτους Πανηγυρικός ]{\textgreek[variant=ancient]{Ἰσοκράτους Πανηγυρικός} 47}

\import{poglavlja/}{pogl08-izokrat-panegirik-47.tex}

\chapter[Ἡροδότου Ἱστορίαι Β]{\textgreek[variant=ancient]{Ἡροδότου Ἱστορίαι Β} 77, 12}

\import{poglavlja/}{pogl09-herodot-2-77.tex}

\chapter[Θουκυδίδου Ἱστορίαι Ϛ 27]{\textgreek[variant=ancient]{Θουκυδίδου Ἱστορίαι Ϛ} 27}

\import{poglavlja/}{pogl10-tukidid-6-27.tex}

\chapter[Λουκιανοῦ Περὶ τοῦ Ἐνυπνίου]{\textgreek[variant=ancient]{Λουκιανοῦ Περὶ τοῦ Ἐνυπνίου ἤτοι Βίος Λουκιανοῦ} 9}

\import{poglavlja/}{pogl20-lukijan-san-9.tex}

\chapter[Ἀντιφῶντος Περὶ ὁμονοίας]{\textgreek[variant=ancient]{Ἀντιφῶντος Περὶ ὁμονοίας} 14}

\import{poglavlja/}{pogl11-antifont-ulomak-14.tex}

\chapter[Φιλοστράτου Tὰ εἰς τὸν Τυανέα Ἀπολλώνιον Ε]{\textgreek[variant=ancient]{Φλαουίου Φιλοστράτου Tὰ εἰς τὸν Τυανέα Ἀπολλώνιον Ε} 14}

\import{poglavlja/}{pogl12-filostrat-apolonije-5-14.tex}

\chapter[Ἐγχειρίδιον Ἐπικτήτου]{\textgreek[variant=ancient]{Ἐγχειρίδιον Ἐπικτήτου} 1}

\import{poglavlja/}{pogl13-epiktet-prirucnik-1.tex}

\chapter[Ξενοφῶντος Ἀπομνημονεύματα Β]{\textgreek[variant=ancient]{Ξενοφῶντος Ἀπομνημονεύματα Β} 4, 5}

\import{poglavlja/}{pogl14-ksen-memor-2-4-5.tex}

\chapter[ Ἀχιλλέως Τατίου Tὰ κατὰ Λευκίππην καὶ Κλειτοφῶντα]{\textgreek[variant=ancient]{ Ἀχιλλέως Τατίου Tὰ κατὰ Λευκίππην καὶ Κλειτοφῶντα Α} 41}

\import{poglavlja/}{pogl15-ahilejt-6-19.tex}

\chapter[Λυσίου Ὀλυμπιακός]{\textgreek[variant=ancient]{Λυσίου Ὀλυμπιακός} 1}

\import{poglavlja/}{pogl16-lizija-olimp-1.tex}

\chapter[Μάρκου Ἀντωνίνου Τῶν εἰς ἑαυτὸν Β]{\textgreek[variant=ancient]{Μάρκου Ἀντωνίνου αὐτοκράτορος \\Τῶν εἰς ἑαυτὸν βιβλίον Β} 11}

\import{poglavlja/}{pogl17-markaur-2-11.tex}

\chapter[Πλωτίνου Ἐννεάδες Ϛ]{\textgreek[variant=ancient]{Πλωτίνου Ἐννεάδες Ϛ} 7, 34}

\import{poglavlja/}{pogl18-plotin-eneade-6-7-34.tex}

\chapter[Μάρκου Ἀντωνίνου Τῶν εἰς ἑαυτὸν Γ]{\textgreek[variant=ancient]{Μάρκου Ἀντωνίνου αὐτοκράτορος \\Τῶν εἰς ἑαυτὸν βιβλίον Γ} 1}

\import{poglavlja/}{pogl19-markaur-3-1.tex}


\end{document}
\clearpage





% kraj
